\chapter{Directivas de aeronavegabilidad: método y casos}\label{ch:21}
\objetivos{Leer una AD desde su aplicabilidad hasta el registro; utilizar las hojas proporcionadas como índices históricos y separar aeronave, motor, hélice y componente.}
\section{Una AD es una acción reglamentaria específica}
Una AD establece acciones para abordar una condición insegura en productos o elementos definidos dentro del sistema de la autoridad emisora. En el sistema europeo, 21.A.3B trata directivas de aeronavegabilidad y M.A.303 o ML.A.303 su cumplimiento aplicable. En Estados Unidos, 14 CFR Parte 39 establece el marco de las AD FAA.

La AD tiene identificación, fecha de efectividad, población, condición insegura, acciones, plazos, referencias y disposiciones adicionales. Se lee completa. El resumen de una base de datos puede orientar la búsqueda, pero no contiene necesariamente las condiciones que determinan si debe ejecutarse una acción.
\section{Aplicabilidad antes que calendario}
Se comprueban fabricante y modelo exactos, números de parte y serie, modificación, instalación y cualquier condición del apartado de aplicabilidad. Una aeronave del mismo modelo puede quedar fuera si no incorpora el equipo o configuración definidos. Una AD de motor puede afectar a aeronaves diferentes que comparten ese producto; por eso un índice de célula no sustituye una búsqueda de motor y componentes.

La clasificación académica separa aeronave, motor, hélice y componente. Una directiva sobre un sistema instalado puede estar publicada bajo varios modelos de aeronave; eso no la convierte en una directiva sobre sus motores. Se conserva el objeto y el nivel de aplicación del texto original.
\section{Plazos y eventos de referencia}
Un plazo puede contarse desde efectividad, desde una utilización acumulada o desde una inspección previa. Puede combinar calendario y horas, o contener un umbral y repetición. «Dentro de 50 horas» sin indicar desde cuándo es una transcripción incompleta. «Salvo que ya se haya realizado» requiere comprobar que la acción previa y su revisión satisfacen la condición de la AD.

Un cumplimiento único no requiere automáticamente repetición; una inspección repetitiva no se cierra para siempre con la primera anotación. Una acción terminal solo elimina repeticiones en las condiciones descritas. Los métodos alternativos de cumplimiento (AMOC) requieren el proceso y autoridad establecidos: una equivalencia técnica propuesta no es un AMOC aprobado.
\section{Qué contienen las ocho hojas aportadas}
Se abrieron las ocho hojas de cálculo de Cessna 172 RG, 172P, 177RG, Piper PA-28RT-201 y 201T, y TB 10, 20 y 200. Contienen, respectivamente, 39, 37, 27, 22, 21, 21, 27 y 12 filas de datos además de cabecera. Hay registros repetidos entre modelos; las 206 filas no representan 206 directivas distintas ni una búsqueda exhaustiva actual.

Las columnas incluyen identificador, asunto, estado indicado en la exportación, fecha de efectividad, referencia, modelo y relaciones de sustitución. El texto «Current» es un dato de la exportación, no prueba automática de vigencia el 18 de octubre de 2026. Se conserva el origen de la hoja y se vuelve al documento oficial.
\section{Caso documental: FAA AD 2024-14-03}
Esta directiva aparece en varias de las hojas. Se contrastó su publicación original en el \textit{Federal Register}, volumen 89, 16 de julio de 2024, pp. 57721--57725, documento 2024-15529. Identificación: AD 2024-14-03, enmienda 39-22784, expediente FAA-2023-1990. El apartado (a) fija efectividad el 20 de agosto de 2024.

El apartado (c) combina modelos de aeronave de su tabla con instalación de un sistema Garmin GFC 500 bajo las condiciones de STC señaladas. La hoja de TB 20 o Cessna no demuestra por sí sola que el equipo esté instalado. El caso ilustra una directiva relacionada con piloto automático, no con el motor.

El procedimiento académico consiste en localizar el registro, abrir el texto oficial, leer (c), contrastar configuración y luego leer (g) y los documentos incorporados por referencia. Se registran acciones, cumplimiento previo admisible y AMOC en los términos del texto. La publicación original confirma esos datos históricos; \textbf{la revisión completa de sustitución/cancelación vigente en DRS queda pendiente y no se declara actualizada la lista de AD de una aeronave real}.
\caso{M21. «No aplicable» necesita motivo}{Una aeronave ficticia coincide con un modelo de la tabla, pero el expediente no contiene información sobre el sistema instalado. La conclusión no es «aplicable» ni «no aplicable» por intuición: es \textit{aplicabilidad pendiente de configuración}. Se solicita identificación del equipo y la modificación. Si se acredita ausencia de la instalación definida, se documenta ese motivo con evidencia; no basta escribir N/A sin referencia.}
\section{Registro de cumplimiento}
La ficha académica contiene AD y revisión, fuente/fecha, producto y configuración, análisis de aplicabilidad, acción y párrafo, documento y revisión usados, fecha y acumulados, resultado, evidencia y próximo vencimiento si existe. Se distingue «aplicable y cumplida», «aplicable pendiente», «no aplicable con evidencia» y «pendiente de determinar».

Un documento incorporado por referencia se conserva con su revisión exacta. La publicación de un SB nuevo no autoriza a cambiar automáticamente la referencia exigida. Si existe AMOC, se registra número, alcance y condiciones y se comprueba que ampare la situación concreta.
\ilustra{flujo-ad}{Secuencia académica de análisis de una directiva. Una respuesta pendiente conserva la tarea abierta hasta obtener evidencia suficiente.}

\section{Una ficha de AD reproducible}
La ficha recoge autoridad, identificador, versión, producto, configuración, efectividad, razón, acción, plazo, repetición, referencias incorporadas, crédito y AMOC cuando existan. Un modelo coincidente puede ser insuficiente si hay condiciones adicionales de instalación o serie. La conclusión enlaza con evidencia de configuración.

Una AD antigua puede seguir siendo pertinente, haber sido sustituida o contener repetición. «Current» en una exportación histórica describe esa exportación. La obra no declara una lista exhaustiva actual de los modelos de las hojas aportadas.
\section{Alcance real de la publicación Garmin}
El original de FAA AD 2024-14-03 verifica identidad histórica, fecha efectiva de 20 de agosto de 2024 y aplicabilidad que combina modelos e instalación definida de Garmin GFC 500. Reconocer TB20 u otro modelo en una fila no completa el análisis: se necesita equipo y configuración.

El original es caso de lectura y no prueba cerrada de vigencia en octubre de 2026. Falta confirmar sustitución, cancelación o modificación en base vigente y resolver aplicabilidad jurisdiccional. Esas casillas permanecen pendientes. La directiva trata piloto automático; no se convierte en una AD del motor.
\section{Disparadores, repetición y crédito}
El plazo puede contar desde efectividad, acumulación total o acción previa; esas palabras son parte de la ecuación. Si la repetición cuenta desde cada inspección, adelantar ejecución puede cambiar el siguiente vencimiento bajo esa regla. Con otra referencia el cálculo difiere. S03 usa una directiva enteramente ficticia para aislar esta operación.

El crédito por trabajo anterior puede exigir acción, revisión y método específicos. Un SB nuevo no modifica por sí solo una AD que incorpora otra revisión. AMOC es una vía regulada, sin libertad para escoger unilateralmente cualquier método supuestamente equivalente.
\repaso{¿La ausencia de equipo en una hoja incompleta prueba no aplicabilidad?}{No demuestra que no esté instalado. Se requiere evidencia de configuración; falta de dato y prueba de ausencia son conceptos distintos.}

\resumen{La búsqueda abarca producto y componentes, además de célula. La hoja es índice, no texto reglamentario. La aplicabilidad precede al cálculo de vencimientos. N/A requiere evidencia; AMOC requiere aprobación en su sistema; un cumplimiento repetitivo conserva un próximo evento.}
\fuentes{Reglamento 748/2012, 21.A.3B \cite{EU748}; 1321/2014, M.A.303 y ML.A.303 \cite{EU1321}; publicación original FAA AD 2024-14-03 \cite{AD2024}; ocho exportaciones aportadas, tratadas como fuentes históricas.}
